\title{xLSTM: Extended Long Short-Term Memory}

\begin{abstract}
During the 1990s, the Long Short-Term Memory (LSTM) model emerged, centered around the innovative ideas of the constant error carousel and gating mechanisms. Over subsequent decades, LSTMs have proven to be both enduring and foundational in a variety of deep learning breakthroughs, notably serving as the basis for the earliest Large Language Models (LLMs). Yet, the introduction of Transformer architectures, grounded in highly parallelizable self-attention, ushered in a new paradigm, quickly surpassing LSTMs at larger scales. This work revisits a straightforward inquiry: To what extent can scaled-up LSTMs—extended to billions of parameters and augmented by contemporary LLM advances—achieve competitive results in language modeling, provided their historical shortcomings are addressed? Our contributions are twofold. First, we present exponential gating enriched by advanced normalization and stabilization strategies. Second, we revise the LSTM memory configuration, resulting in: (i) the sLSTM, which incorporates a scalar-valued memory, scalar-based updates, and an enhanced memory mixing process, and (ii) the mLSTM, which is amenable to full parallelization by virtue of its matrix memory and covariance-based update rules. By embedding these LSTM variants within residual block frameworks, we formulate xLSTM blocks, which can be recursively assembled into deep xLSTM networks. The synergy of exponential gating and reimagined memory structures substantially elevates the performance of xLSTM models, enabling them to effectively rival leading-edge Transformers and State Space Models in both scalability and results.\\
Code available at: \url{https://github.com/NX-AI/xlstm}
\end{abstract}

\section{Introduction}
\label{sec:intro}

The Long Short-Term Memory (LSTM) ideas~\citep{Hochreiter:91,Hochreiter:97nips,Hochreiter:97}, i.e., the constant error carousel and gating, were introduced to overcome the vanishing gradient problem of recurrent neural networks~\citep{Hochreiter:91,Hochreiter:00book}:

\begin{align}
\label{eq:lstm_idea}
  c_t \ = \ f_t \ c_{t-1} \ + \ 
          i_t \  z_t \ , \quad  
          h_t \ = \ o_t \ \psi({c_t}) \ . 
\end{align}

The constant error carousel refers to the process where the previous cell state $c_{t-1}$ (indicated in green) is incrementally adjusted by the cell inputs $z_t$, with this modification being regulated through sigmoid gates (depicted in blue). The input gate $i_t$ and forget gate $f_t$ are responsible for governing these adjustments, while the output gate $o_t$ determines the information that is emitted from the memory cell, specifically the hidden state $h_t$. Prior to output, the cell state undergoes normalization or squashing via $\psi$, after which the output gating yields the hidden state.

LSTMs have demonstrated substantial success across a wide array of applications \citep{Hochreiter:01,Hochreiter:07,Schmidhuber:15}, dominating the field of text generation until they were largely supplanted by Transformers in 2017~\citep{Vaswani:17}. Their capability has been validated on numerous sequence-based problems, including but not limited to, tasks such as text generation~\citep{Graves:13,Karpathy:15blog}, handwriting synthesis~\citep{Graves:13}, sequence-to-sequence machine translation~\citep{Sutskever:14nips}, program evaluation~\citep{Zaremba:14arxiv}, image captioning~\citep{Karpathy:15,Hossain:19}, code generation~\citep{Karpathy:15blog}, modeling rainfall-runoff processes~\citep{Kratzert:18,Kratzert:19}, and developing hydrological models used for flood warning systems~\citep{Nearing:24}. Within reinforcement learning, LSTMs remain the preeminent architecture among sequence models, as evidenced by their use in the AlphaStar agent for StarCraft II~\citep{Vinyals:17short}, OpenAI Five for Dota 2~\citep{OpenAI:19}, and in control systems for nuclear fusion magnetic controllers~\citep{Degrave:22short}. Notably, LSTMs are particularly proficient in forming abstractions by extracting and retaining semantic features within their memory cells~\citep{Karpathy:15blog}, a strength illustrated through specialized units such as number and syntax neurons~\citep{Lakretz:19}, linguistic neurons~\citep{Bau:19}, and sentiment neurons~\citep{Radford:17arxiv}. LSTMs continue to see deployment in impactful, real-world contexts~\citep{Degrave:22short,Nearing:24}, illustrating their enduring relevance and reliability.

Although LSTMs have achieved significant accomplishments, they exhibit three principal drawbacks: (i) Inability to modify previously made storage decisions. This drawback is illustrated through the {\it Nearest Neighbor Search} problem (refer also to Appendix~\ref{sec:appNearestNeighborSearch}): Provided with a reference vector, the sequence must be processed in order to identify the vector most similar to the reference, thereby retrieving its corresponding value at the end. The left panel of Figure~\ref{fig:lstmProblems} displays the mean squared error in this task. LSTMs encounter difficulties in overwriting a stored value upon discovering a more similar vector, whereas the proposed xLSTM mitigates this issue through exponential gating mechanisms. (ii) Restricted storage capacity, as all relevant information is forced into a scalar cell state. This is exemplified by {\it Rare Token Prediction}. The right panel of Figure~\ref{fig:lstmProblems} presents token prediction perplexity on Wikitext-103~\citep{Merity:17}, segmented according to token frequency. LSTM performance deteriorates for infrequent tokens due to its constrained memory, but this shortcoming is addressed by xLSTM’s implementation of matrix-based memory. (iii) A deficiency in parallelism caused by memory entanglement; specifically, the hidden-to-hidden connections between consecutive time steps necessitate sequential operations.

The aforementioned drawbacks of the LSTM architecture have contributed to the rise of Transformer models~\citep{Vaswani:17} in the language modeling domain. An important question emerges: how effective might language modeling become if these constraints were removed and LSTMs were scaled to the dimensions of current Large Language Models?

\section{Extended Long Short-Term Memory}
\label{sec:xLSTM}

To address the inherent constraints of LSTMs, Extended Long Short-Term Memory (xLSTM) incorporates two significant advancements to the core LSTM concept presented in Equation~\eqref{eq:lstm_idea}. These advancements consist of exponential gating mechanisms and innovative memory structures, giving rise to two new variants within the LSTM family: (i) the sLSTM (refer to Section~\ref{sec:sLSTM}), which employs scalar-valued memory, scalar updates, and inter-cell memory mixing; and (ii) the mLSTM (see Section~\ref{sec:mLSTM}), which utilizes matrix-based memory and leverages a covariance-style (outer product) update mechanism, thereby allowing full parallel computation. Both sLSTM and mLSTM leverage the concept of exponential gating for improved performance. To further support parallel operations, mLSTM forgoes memory mixing—effectively removing recurrent connections among hidden units. Both architectures, sLSTM and mLSTM, permit the inclusion of multiple memory cells; in sLSTM, memory mixing can occur between cells. Additionally, sLSTM allows for a multi-head configuration, in which memory mixing occurs exclusively among cells within a head and not across different heads, introducing a refined form of memory interaction thanks to exponential gating. For mLSTM, deploying multiple heads and multiple cells is functionally interchangeable.

By embedding these LSTM extensions as residual block components, xLSTM blocks are formed (explained in Section~\ref{sec:xLSTMblock}). Stacking such xLSTM blocks residually within the overall model structure defines the xLSTM architectures (see Section~\ref{sec:xLSTMarchitecure}). The structure and constituent parts of the xLSTM architecture are depicted in Figure~\ref{fig:xlstm_sketch}.

\subsection{Review of the Long Short-Term Memory}
\label{sec:orgLSTM}

The foundational concept behind LSTM~\citep{Hochreiter:91,Hochreiter:97nips,Hochreiter:97} featured a scalar memory cell functioning as the principal unit for computation and storage, mitigating the vanishing gradient issue~\citep{Hochreiter:91,Hochreiter:00book} by maintaining stable error flow through the constant error carousel, represented by the cell state's update mechanism. This memory cell architecture encompasses three types of gates: input, output, and forget gates, with the forget gate being subsequently added by \citet{Gers:00}. The LSTM memory cell's update equations at time step $t$ are as follows:

\begin{align}
c_t \ &= \  f_t \ c_{t-1} \ + \ i_t \ z_t &  & &\text{cell state} \\
h_t  \ &= \ o_t \ \tilde{h}_t \ , & \tilde{h}_t \ &= \ \psi \left( c_t\right)
  &\text{hidden state} \\
z_t \ &= \ \varphi \left( \tilde{z}_t \right) \ , 
  &\tilde{z}_t \ &=  \ \mathbf{w}^\top_{z} \ \mathbf{x}_t \ + \
  r_{z}  h_{t-1} \ + \  b_{z} \ \
  &\text{cell input} \\
i_t \ &= \ \sigma \left( \tilde{i}_t  \right) \ , 
  &\tilde{i}_t \ &= \ \mathbf{w}^\top_{i} \ \mathbf{x}_t \ + \
  r_{i}  \ h_{t-1} \ + \  b_{i} \ \
  &\text{input gate} \\
f_t \ &= \ \sigma \left(  \tilde{f}_t \right) \ , 
  &\tilde{f}_t \ &= \ \mathbf{w}^\top_{f} \ \mathbf{x}_t  \ + \
  r_{f}  \ h_{t-1} \ + \  b_{f} \ \
  &\text{forget gate} \\
o_t \ &= \ \sigma \left( \tilde{o}_t \right) \ , 
  &\tilde{o}_t  \ &= \ \mathbf{w}^\top_{o} \ \mathbf{x}_t \ + \
  r_{o}  \ h_{t-1} \ + \  b_{o} \ \
  &\text{output gate} 
\end{align}

The vectors $\mathbf{w}_{z}$, $\mathbf{w}_{i}$, $\mathbf{w}_{f}$, and $\mathbf{w}_{o}$ serve as the weights connecting the input $\mathbf{x}_t$ to the cell input, input gate, forget gate, and output gate, respectively. Similarly, $r_{z}$, $r_{i}$, $r_{f}$, and $r_{o}$ are the recurrent weights from the previous hidden state $h_{t-1}$ to each of these respective components. The biases for each element are denoted as $b_{z}$, $b_{i}$, $b_{f}$, and $b_{o}$. The activation functions for the cell input and hidden state are represented by $\varphi$ and $\psi$, where both are most often the hyperbolic tangent ($\tanh$). The function $\psi$ constrains or compresses the values of the cell state, ensuring it remains within a bounded range. Each gating mechanism applies a sigmoid nonlinearity, defined as $\sigma \left( x \right)= 1/(1+\exp(-x))$. In subsequent advances, several scalar memory cells are organized into a vector $\mathbf{c}_t \in \mathbb{R}^{d}$, permitting the use of a recurrent weight matrix $\mathbf{R} \in \mathbb{R}^{d \times d}$ to blend outputs among memory cells~\citep{Greff:15}; further explanation is provided in Appendix~\ref{sec:apporgLSTM}. Results from ablation experiments indicate that every part of the memory cell architecture plays an essential role~\citep{Greff:15}.

\subsection{sLSTM}
\label{sec:sLSTM}

To enable LSTMs to modify their storage choices, we propose the use of exponential gates (shown in red), complemented by normalization and stabilization techniques. Specifically, the input and forget gates are augmented with exponential activation functions. For normalization purposes, we add a normalizer state, which accumulates the sum of the input gate multiplied by all subsequent forget gates.

The scalar sLSTM forward pass is:

\begin{align}
\label{eq:slstmforward}
c_t \ &= \  f_t \ c_{t-1} \ + \ i_t \ z_t & & 
 &\text{cell state} \\
n_t \ &= \  f_t \ n_{t-1} \ + \ i_t  & & 
 &\text{normalizer state} \\
h_t \ &= \ o_t \ \tilde{h}_t \ ,
  &\tilde{h}_t \ &= \ c_t / n_t
&\text{hidden state} \\
z_t \ &= \ \varphi \left( \tilde{z}_t \right) \ , 
  &\tilde{z}_t \ &=  \ \mathbf{w}^\top_{z} \ \mathbf{x}_t \ + \
  r_{z}  h_{t-1} \ + \  b_{z}
  &\text{cell input} \\
i_t \ &= \ \!\exp \left( \tilde{i}_t  \right) \ , 
  &\tilde{i}_t \ &= \ \mathbf{w}^\top_{i} \ \mathbf{x}_t \ + \
  r_{i}  \ h_{t-1} \ + \  b_{i}
  &\text{input gate} \\
f_t \ &= \ \sigma \left(  \tilde{f}_t \right) \ \text{OR} \
\!\exp \left(  \tilde{f}_t \right) \ ,
  &\tilde{f}_t \ &= \ \mathbf{w}^\top_{f} \ \mathbf{x}_t  \ + \
  r_{f}  \ h_{t-1} \ + \  b_{f}
  &\text{forget gate} \\
o_t \ &= \ \sigma \left( \tilde{o}_t \right) \ , 
  &\tilde{o}_t  \ &= \ \mathbf{w}^\top_{o} \ \mathbf{x}_t \ + \
  r_{o}  \ h_{t-1} \ + \  b_{o}
  &\text{output gate} 
\end{align}

We transfer the original LSTM gating techniques, i.e., input- and/or hidden-dependent gating plus bias term, to the new architectures. Exponential activation functions can lead to large values that cause overflows. Therefore, we stabilize gates with an additional state \mbox{$m_t$~\citep{Milakov:18arxiv}}:

\begin{align}
\label{eq:slstmstabil}
m_t \ &= \ \max \left( \log ( f_t ) + m_{t-1} , \log ( i_t ) \right) &\text{stabilizer state} \\
i'_t \ &= \ \!\exp \left( \log \left ( i_t \right) - m_t \right) \ = \ \!\exp \left( \tilde{i}_t - m_t \right) \ \, 
  &\text{stabil. input gate} \\
f'_t \ &= \ \!\exp \left( \log \left( f_t \right) + m_{t-1} - m_t \right) \ \, 
  &\text{stabil. forget gate}
\end{align}

Appendix~\ref{sec:appsLSTM} demonstrates that substituting $f_t$ with $f'_t$ and $i_t$ with $i'_t$ during the forward computation leaves both the network’s outputs and the gradients of the loss with respect to its parameters unaltered.

\paragraph{New Memory Mixing.} As with the standard LSTM, sLSTM allows for several memory cells (refer to Appendix~\ref{sec:appsLSTM} for details). This configuration facilitates memory mixing, achieved through recurrent links $\mathbf{R}_{\mathbf{z}}$, $\mathbf{R}_{\mathbf{i}}$, $\mathbf{R}_{\mathbf{f}}$, $\mathbf{R}_{\mathbf{o}}$ that route signals from the hidden state $\mathbf{h}$ to the input of the memory cell $\mathbf{z}$, as well as to the gating vectors $\mathbf{i}$, $\mathbf{f}$, $\mathbf{o}$. Exponential gating introduces a novel dimension to memory mixing. The updated sLSTM architecture permits multiple heads, enabling intra-head memory mixing while preventing inter-head interaction. The combination of these heads and exponential gating yields an innovative mechanism for memory mixing within the sLSTM framework.

\subsection{mLSTM}
\label{sec:mLSTM}

In order to bolster the memory capabilities of LSTMs, we replace the conventional scalar memory cell $c \in \mathbb{R}$ with a matrix representation $\mathbf{C} \in \mathbb{R}^{d \times d}$. This allows information retrieval through matrix multiplication operations. At any time step $t$, a pair consisting of a key vector $\mathbf{k}_t \in \mathbb{R}^d$ and a value vector $\mathbf{v}_t \in \mathbb{R}^d$—using the terminology from the Transformer architecture—is to be stored. Subsequently, at a future time step $t+\tau$, the stored value $\mathbf{v}_t$ should be accessible via a query vector $\mathbf{q}_{t+\tau} \in \mathbb{R}^d$. This memory management scenario is consistent with the framework of Bidirectional Associative Memories (BAMs) as described in \citep{Kohonen:72,Anderson:72,Nakano:72,Anderson:77}.

The covariance update rule~\citep{Sejnowski:77,Dayan:91} for storing a key-value pair is

\begin{align}
\mathbf{C}_t \ &= \ \mathbf{C}_{t-1} \ + \ \mathbf{v}_{t} \ \mathbf{k}_t^\top \ .
\end{align}

We posit that a layer normalization is applied prior to projecting the inputs to key and value vectors, thus ensuring a zero mean. The covariance-based update mechanism achieves optimality~\citep{Dayan:91} in maximizing the separability of recovered binary vectors, effectively corresponding to an optimal signal-to-noise ratio. If retrieval is constrained to involve only pairwise interactions, even greater separability can be achieved, albeit with the increase to quadratic computational complexity, as seen in attention mechanisms~\citep{Krotov:16,Krotov:17,Ramsauer:21}. This same update approach matches the Fast Weight Programmers~\citep{Schmidhuber:92ncfastweights,Schlag:21}, where more recent versions introduce a constant decay rate that modulates $\mathbf{C}_{t-1}$ and apply a constant learning rate to 
$\mathbf{v}_{t} \mathbf{k}_t ^\top$~\citep{Ba:16}. Following this rationale, we incorporate the covariance-based update within the LSTM architecture, identifying the forget gate with the decay rate, the input gate with the learning rate, and utilizing the output gate to modulate the recovered vector.

For this matrix memory, the normalizer state is the weighted sum of key vectors, where each key vector is weighted by the input gate and all future forget gates. Again, the normalizer state keeps record of the strength of the gates. Since the dot product between query and normalizer state can be close to zero, we use the absolute value of this dot product and lower bound it by a threshold (typically 1.0) as done previously~\citep{Sun:23arxiv}. The mLSTM forward pass is:

\begin{align}
\label{eq:mlstm_recurrent_begin}
\mathbf{C}_t \ &= \  f_t \ \mathbf{C}_{t-1} \ + \ 
  i_t \ \mathbf{v}_t \ \mathbf{k}_t^\top & &  &\text{cell state} \\
\mathbf{n}_t \ &= \  f_t \ \mathbf{n}_{t-1} \ + \ 
  i_t \ \mathbf{k}_t & &  &\text{normalizer state} \\
\mathbf{h}_t  \ &= \ \mathbf{o}_t \ \odot \ \tilde{\mathbf{h}}_t
\ , \qquad \qquad 
\tilde{\mathbf{h}}_t \ = \   \mathbf{C}_t \mathbf{q}_t \ / \ 
  \max \left\{ |\mathbf{n}_t^\top \mathbf{q}_t|, 1 \right\} 
   &&&\text{hidden state} \\
\mathbf{q}_t \ &= \ \mathbf{W}_q \ \mathbf{x}_t \ + \ \mathbf{b}_q  & &  &\text{query input} \\
\mathbf{k}_t \ &= \ \frac{1}{\sqrt{d}} \mathbf{W}_k \ \mathbf{x}_t \ + \ \mathbf{b}_k  & &  &\text{key input} \\
\mathbf{v}_t \ &= \ \mathbf{W}_v \ \mathbf{x}_t \ + \ \mathbf{b}_v  & &  &\text{value input} \\
i_t \ &= \ \!\exp \!  \! \left( \tilde{i}_t  \right) \ , 
 \qquad \qquad \ \ \ \ \,
  \tilde{i}_t \ = \ \mathbf{w}^\top_{i} \ \mathbf{x}_t \ + \  b_{i} \ \ \ 
  &&&\text{input gate} \\
f_t \ &= \ \sigma \!  \left(  \tilde{f}_t \right) \ \text{OR} \ \!\exp \!  \! \left(  \tilde{f}_t \right) , \ \ \: \!
\tilde{f}_t \ = \ \mathbf{w}^\top_{f} \ \mathbf{x}_t  \ + \
  b_{f}
  &&& \text{forget gate} \\
\label{eq:mlstm_recurrent_end}
\mathbf{o}_t \ &= \ \sigma \left( \tilde{\mathbf{o}}_t \right) \ , \qquad \qquad \quad \ \ \ \,
  \tilde{\mathbf{o}}_t  \ = \ \mathbf{W}_{\mathbf{o}} \ \mathbf{x}_t \ + \
  \mathbf{b}_{\mathbf{o}}
  &&&\text{output gate} 
\end{align}

mLSTM, similar to standard LSTM architectures, supports the use of several memory cells. For mLSTM, having multiple heads and multiple cells is functionally indistinguishable due to the absence of any mechanism for memory intermixing. To prevent instability arising from the exponential gating in mLSTM, we apply the stabilization strategies utilized for sLSTM, as described in Equation~\ref{eq:slstmstabil}. Given that mLSTM does not incorporate memory mixing, its recurrence relation can be recast into a parallel formulation. Additional information is provided in Appendix~\ref{sec:appmLSTM}.

\subsection{xLSTM Architecture}
\label{sec:xLSTMarch}

\paragraph{xLSTM Blocks.}
\label{sec:xLSTMblock}
The xLSTM block is designed to encapsulate prior information through non-linear transformation within a high-dimensional embedding, thereby facilitating a clearer distinction among diverse historical states or contexts. This separation is fundamental for accurate next-element prediction in sequences (e.g., predicting the following token). This approach draws inspiration from Cover’s Theorem~\citep{Cover:65}, which posits that patterns transformed non-linearly into higher dimensional spaces are statistically more likely to become linearly separable compared to their representation in lower dimensions. Two primary residual block designs are investigated: (i) A residual block featuring post up-projection—similar to the Transformer structure—which first generates a non-linear summary of previous inputs in the original representation space, then performs a linear projection into an expanded dimension, applies a non-linear activation, and projects the output back into the original space; this is illustrated in the left section of Figure~\ref{fig:Backbones}
and in the third column of Figure~\ref{fig:xlstm_sketch}. A further breakdown is available in Figure~\ref{fig:appsLSTM_detailed} in the appendix. (ii) An alternative is the residual block with pre up-projection (analogous to those in State Space Models), in which a linear mapping to a higher-dimensional space is carried out first,

encodes previous information non-linearly within a higher-dimensional representation, which is subsequently mapped back to the input space through a linear transformation. When an xLSTM block utilizes an sLSTM module, we apply the up-projection after this block. In contrast, for xLSTM blocks that employ an mLSTM, the up-projection occurs beforehand, as the increased dimensionality expands the available memory capacity. For a visual explanation, consult the leftmost panel of Figure~\ref{fig:Backbones}, the third column of Figure~\ref{fig:xlstm_sketch}, or Figure~\ref{fig:appsLSTM_detailed} in the appendix.

\paragraph{xLSTM Architecture.}
\label{sec:xLSTMarchitecure}
The xLSTM architecture is built by stacking modular units in a residual fashion, following the design principles in~\citep{Srivastava:15,He:16}. Our implementation adopts a pre-LayerNorm~\citep{Ba:16b} residual backbone, which is the prevailing architecture in recent Large Language Models. Refer to the final column of Figure~\ref{fig:xlstm_sketch} for further illustration.

\subsection{Memory and Speed Considerations}

In contrast to Transformers, xLSTM architectures exhibit linear computational demands and constant memory usage relative to sequence length. The compressive nature of xLSTM memory makes it particularly advantageous for deployment in industrial scenarios and at the edge.

The mLSTM’s memory does not depend on learnable parameters, yet it incurs computational cost due to the $d \times d$ memory matrix and its corresponding $d \times d$ updates. This results in a trade-off between available memory capacity and computational overhead. Despite this, the operations are amenable to parallel execution on GPUs, which minimizes their impact on overall runtime.

Although mLSTM can be parallelized in a fashion similar to FlashAttention~\citep{Dao:22,Dao:23} or GLA~\citep{Yang:23arxiv}, sLSTM cannot exploit such parallelism owing to its use of memory mixing via hidden-to-hidden connections. Nevertheless, by implementing a highly optimized CUDA version that leverages GPU registers for memory access, we achieve sLSTM speeds that are usually less than twice as slow as those for mLSTM.

\section{Related Work}
\label{sec:related}

\paragraph{Linear Attention.} Numerous approaches have been proposed to address the quadratic scaling of Transformer attention with respect to sequence length, aiming to achieve linear complexity. The Synthesizer generates attention weights synthetically, sidestepping direct token-to-token interactions~\citep{Tay:20arxiv}. Linformer employs a low-rank matrix to implement self-attention and further introduces linear approximations~\citep{Wang:20arxiv}. Linear Transformer modifies the attention process to allow for linearization~\citep{Katharopoulos:20}. Performer utilizes a positive orthogonal random features method to produce a linear approximation of the attention softmax~\citep{Choromanski:21}. Alternatively, traditional attention mechanisms have been replaced with efficient long convolutions, as exemplified by Structured Global Convolution (SGConv)~\citep{Li:22} and Hyena Hierarchy~\citep{Poli:23}.

\paragraph{State Space Models.} State Space Models (SSMs) have recently attracted substantial attention due to their linear computational complexity with respect to input sequence length and their competitive results against Transformer architectures. Early milestones in this domain include the Structured State Space sequence model (S4)~\citep{Gu:21}, with subsequent developments such as the Diagonal State Space (DSS) model~\citep{Gupta:22}, Gated State Space (GSS) models~\citep{Mehta:22}, S5 model~\citep{Smith:22}, Bidirectional Gated SSM (BiGS)~\citep{Wang:22}, H3 model~\citep{Fu:23}, and Mamba~\citep{Gu:24arxiv}.

\paragraph{Recurrent Neural Networks.} In recent work, Recurrent Neural Networks (RNNs) have been proposed as alternatives to Transformer and attention mechanisms, primarily because their computational complexity scales linearly with sequence length. Language modeling has benefited from variants like Deep Linear Recurrent Units (LRUs)~\citep{Orvieto:23, De:24arxiv}, along with the Hierarchically Gated Linear RNN (HGRN)~\citep{Qin:23} and its successor HGRN2~\citep{Qin:24arxiv}, which have both delivered favorable outcomes. Among established RNN methods for large-scale language modeling, RWKV~\citep{Peng:23arxivshort,Peng:24arxivshort} stands out, demonstrating performance levels that rival Transformer architectures.

\paragraph{Gating.}
A central innovation within LSTM architectures is the introduction of gating mechanisms—a concept that has been independently revisited and reconceptualized by numerous contemporary methods. This gating paradigm appears in works such as HGRN~\citep{Qin:23}, HGRN2~\citep{Qin:24arxiv}, Gated Linear Attention (GLA)~\citep{Yang:23arxiv}, Gated State Space (GSS) models~\citep{Mehta:22}, Bidirectional Gated SSM (BiGS)~\citep{Wang:22}, Moving Average Equipped Gated Attention (MEGA)~\citep{Ma:22}, as well as more recent architectures like RWKV~\citep{Peng:23arxivshort} and Mamba~\citep{Gu:24arxiv}.

\paragraph{Covariance Update Rule.}
In an effort to improve memory capacity, the mLSTM cell was augmented with a matrix-based memory module employing a covariance update rule. Comparable strategies leveraging covariance-style updates have been employed by models such as Fast Weight Programmers \citep{Schmidhuber:92ncfastweights,Schlag:21}, RWKV-5 and RWKV-6~\citep{Peng:24arxivshort}, Retention~\citep{Sun:23arxiv}, Linear Transformer~\citep{Katharopoulos:20}, along with HGRN2~\citep{Qin:24arxiv}.

\paragraph{Most Related.} The models most akin to xLSTM in terms of architecture are Retention~\citep{Sun:23arxiv}, RWKV~\citep{Peng:23arxivshort,Peng:24arxivshort}, and HGRN2~\citep{Qin:24arxiv}. These frameworks incorporate mechanisms such as matrix memory and/or gating. Nevertheless, unlike the newly proposed sLSTM, these alternatives do not accommodate memory mixing. The capacity for memory mixing is critical for addressing state tracking tasks, positioning LSTMs as more powerful than both State Space Models (SSMs) and Transformers~\citep{Merrill:24,Deletang:23}. State tracking is essential for activities such as code execution or monitoring entities throughout extended narratives.

\paragraph{Residually Stacking Architectures.} Modern large-scale deep learning models, including xLSTM, are typically designed by sequentially stacking components with residual connections~\citep{Srivastava:15,He:16}.

This residual stacking methodology facilitated the development of deep convolutional architectures~\citep{He:16} and paved the way for the rise of Transformers~\citep{Vaswani:17}. Transformers form the foundational backbone for Large Language Models (LLMs) such as GPT-3~\citep{Brown:20short}, ChatGPT~\citep{Schulman:22short}, GPT-4~\citep{Achiam:23gpt4short}, Megatron-LM~\citep{Shoeybi:19}, Gopher~\citep{Rae:21short}, ERNIE 3.0 Titan~\citep{Wang:21arxivshort}, GLaM~\citep{Du:21glamshort}, Chinese M6~\citep{Lin:21}, multilingual AlexaTM 20B~\citep{Soltan:22}, OPT~\citep{Zhang:22opt}, Chinchilla~\citep{Hoffmann:22short}, BLOOM~\citep{Scao:22arxivshort}, GLM-130B~\citep{Zeng:22arxivshort}, LaMDA~\citep{Thoppilan:22short}, PaLM~\citep{Chowdhery:22short}, Llama~\citep{Touvron:23arxiv}, Gemini~\citep{GeminiTeam:23arxiv,Reid:24arxiv}.

\section{Experiments}
\label{sec:Exp}

In this section, we present an empirical analysis of xLSTM, positioning its performance against other established approaches in the context of language modeling. Section~\ref{sec:ExpSyntethic} is dedicated to examining xLSTM's unique abilities through a series of synthetic tasks. In Section~\ref{sec:ExpComparison}, we report the validation perplexity achieved by several modern language modeling techniques, each trained on a corpus of 15B tokens from SlimPajama~\citep{Soboleva:23}, including a suite of ablation experiments focusing on xLSTM itself. Furthermore, we explore the scaling trends exhibited by each model in relation to findings from \citet{Kaplan:20} and \citet{Brown:20short}.

Section~\ref{sec:ExpLanguage} delves into more comprehensive experiments on language modeling. Here, xLSTM and the leading methods identified in Section~\ref{sec:ExpComparison} are compared after training on an expanded dataset consisting of 300B tokens from SlimPajama~\citep{Soboleva:23}. Our analysis proceeds in four steps: firstly, we measure how each approach generalizes to longer input contexts; secondly, we evaluate their performance by validation perplexity and a range of downstream tasks~\citep{Sutawika:24}; thirdly, we test them across 571 textual domains included in the PALOMA language benchmark dataset~\citep{Magnusson:23arxivshort}; and fourthly, we revisit scaling properties with a significantly larger training set, representing a 20-fold increase in data.

\label{sec:xlstm_blocks_notation}
Throughout all experiments, we denote the ratio of mLSTM- to sLSTM-based xLSTM blocks by xLSTM[$a$:$b$], where $a/b$ reflects this proportion. As an illustration, xLSTM[7:1] indicates that seven out of eight blocks employ mLSTM, while one utilizes sLSTM. In cases where the total number of blocks is 48, this configuration corresponds to 42 blocks based on mLSTM and 6 on sLSTM. Additionally, in all experiments, up-projection blocks are applied before and after mLSTM and sLSTM blocks, respectively.

\subsection{Synthetic Tasks and Long Range Arena}
\label{sec:ExpSyntethic}
To begin, we evaluate how well xLSTM's novel exponential gating mechanism with memory mixing operates on formal language benchmarks~\citep{Deletang:23}. Subsequently, we analyze the impact of xLSTM's matrix memory component on the Multi-Query Associative Recall task~\citep{Arora:23arxiv}. Lastly, we measure the capability of xLSTM to handle long-context sequences using the Long Range Arena suite~\citep{Tay:21}.

\paragraph{Evaluating xLSTM’s Exponential Gating with Memory Mixing.} We assess the impact of xLSTM’s recently introduced exponential gating mechanism, which incorporates memory mixing, purported to enhance its capability for state tracking tasks~\citep{Merrill:24, Merrill:23}. To this end, we adapt and extend the suite of formal language tasks presented in \citet{Deletang:23}, enabling model training on sequences of multiple lengths to study extrapolation behaviors. Comprehensive task definitions and further experimental details are provided in Appendix~\ref{sec:appExpSynthetic-formalLang}. 

Our study benchmarks xLSTM against several architectures, including Transformers, State Space Models, and traditional Recurrent Neural Networks, by computing accuracy metrics for tokens pertinent to each task. The accuracy metric ranges continuously from 0 (equivalent to chance) to 1 (error-free performance). Specifically, we experiment with 2-block variants of each model: xLSTM[0:1] (employing sLSTM only), xLSTM[1:0] (using mLSTM only), xLSTM[1:1], as well as Llama, Mamba, RWKV, Retention, Hyena, LSTM, and LSTM implemented within Transformer architectures (LSTM (Block)). Results are visualized in Figure~\ref{fig:formal-main}. 

We find that architectures lacking a memory mixing mechanism—such as plain Transformers or State Space Models—are unable to tackle certain regular grammar tasks (for example, the parity task), reflecting their inability to maintain effective state tracking. These results align with prior evidence suggesting Transformers and State Space Models are intrinsically less capable than RNNs in this domain~\citep{Merrill:24, Merrill:23, Deletang:23}.

\paragraph{Evaluation of xLSTM's Memory on Associative Recall Tasks.} This study assesses the memory capacity of xLSTM's newly introduced matrix memory using the Multi-Query Associative Recall task~\citep{Arora:23arxiv}. In this setup, each input sequence contains key-value pairs selected randomly from an extensive vocabulary, and the challenge is to store these associations for subsequent retrieval. To make the task more challenging than in its original design, the number of key-value pairs is increased to as many as 256, and the sequence length can reach up to 2048 tokens, thereby providing a more rigorous assessment of model memory. We conduct comparisons among various two-block model designs, including Llama, Mamba, RWKV-5, RWKV-6, xLSTM[1:1], and xLSTM[1:0]. The evaluation metric is the accuracy of successfully recalling key-value pairs. Because the Transformer architecture (e.g., Llama) exhibits memory capacity that scales exponentially with the coding dimension~\citep{Ramsauer:21}, it sets the benchmark for this evaluation. Figure~\ref{fig:mqar-main} displays the corresponding results. Notably, xLSTM[1:1] achieves the strongest performance among all tested non-Transformer models, a trend observable even in smaller model sizes. Counter to intuition, including the sLSTM block does not restrict memory performance; instead, it contributes positively, especially noticeable at the highest task difficulty involving 256 key-value pairs. Supplementary findings are available in Appendix~\ref{sec:appExpSynthetic-mqar}, where further analysis demonstrates that the improved memory of xLSTM enables generalization to sequence lengths exceeding those encountered during training.

\paragraph{Test of xLSTM's Long Context Capabilities on Long Range Arena.} In evaluating xLSTM's effectiveness on lengthy sequences and extended contexts, we benchmark various approaches using the Long Range Arena~\citep{Tay:21}. Across all tasks, xLSTM yields robust and stable results, indicating that its architecture is highly capable of efficiently addressing diverse challenges related to long context processing.

For more details, see Appendix~\ref{sec:appExpSynthetic-lra}.

\subsection{Method Comparison and Ablation Study}
\label{sec:ExpComparison}

In order to investigate the primary research question—specifically, the capabilities of our novel LSTM variants when extensively scaled for language modeling—we conduct training of xLSTMs alongside Transformers, State Space Models, and additional baselines. All models are trained on 15B tokens drawn from the SlimPajama dataset within a consistent auto-regressive framework. After training, we assess their performance on the validation set and carry out ablation experiments on the xLSTMs to further analyze their behavior.

\paragraph{Comparing xLSTM to Other Methods.} To facilitate comparison, all models were trained on a dataset containing 15B tokens from SlimPajama~\citep{Soboleva:23}. Their performance was assessed through perplexity measured on a held-out validation set. We evaluate a range of approaches: xLSTM (our proposed method), GPT-3 (Transformer)~\citep{Brown:20short}, Llama (Transformer)~\citep{Touvron:23arxiv}, H3 (SSM)~\citep{Fu:23}, Mamba (SSM)~\citep{Gu:24arxiv}, RWKV-4 (RNN)~\citep{Peng:23arxivshort}, RWKV-5 (RNN)~\citep{Peng:24arxivshort}, RWKV-6 (RNN)~\citep{Peng:24arxivshort}, GLA (linear Transformer)~\citep{Yang:23arxiv}, HGRN (RNN)~\citep{Qin:23}, HGRN2 (RNN)~\citep{Qin:24arxiv}, RetNet (linear Transformer)~\citep{Sun:23arxiv}, Hyena (linear Transformer)~\citep{Poli:23}, as well as variants xLSTM[1:0] and xLSTM[7:1] as described in Section~\ref{sec:xlstm_blocks_notation}. Model training was conducted using mixed precision where possible. Notably, for RWKV-5, RWKV-6, GLA, and HGRN2, mixed precision training was achieved without relying on the automated mixed precision provided by PyTorch (refer also to Appendix Section~\ref{sec:appGenTrainProc} for further explanation).

The evaluated techniques fall under three principal categories: (a) Transformers, (b) State Space Models (SSMs), and (c) Recurrent Neural Networks (RNNs) and linear Transformers. Linear Transformers are characterized by linear-time alternatives to the canonical Transformer attention. All models in the comparison are constructed to match the parameter count of a 350M parameter GPT-3 configuration, specifically with an embedding dimension of 1024 and comprising 24 residual blocks. It should be noted that only GPT-3 applies weight sharing between input and output token embeddings, resulting in a slightly reduced total parameter count.

As indicated in Table~\ref{tab:spaj15b_model_results}, the xLSTM achieves superior validation perplexity compared to all other models under consideration. Additional implementation specifics are supplied in Appendix~\ref{sec:appExpComparison}. Furthermore, Figure~\ref{fig:temp_sclaw15B} presents scaling trends from this experiment, highlighting that xLSTM's advantages are expected to persist as model size increases.

\paragraph{Ablation Studies.}
As shown in Table~\ref{tab:spaj15b_model_results} and Figure~\ref{fig:temp_sclaw15B}, the xLSTM architecture attains superior performance in language modeling when trained with 15B tokens from SlimPajama. To isolate the impact of each modification leading from a standard LSTM to xLSTM, we incrementally adapt the vanilla LSTM design to align with xLSTM's structure. The process begins by incorporating LSTM layers into residual backbones featuring pre-LayerNorm. Subsequently, this is augmented with a post up-projection block. The final transformation introduces exponential gating mechanisms and matrix memory components.

The results are shown in Table~\ref{tab:ablstudies} (top). 

The ablation experiments credit the substantial boost in performance to a combination of the exponential gating and the matrix memory. Furthermore, recognizing the critical role of gating within RNNs and State Space Models, we investigate the effects of various gating strategies. As shown in Table~\ref{tab:ablstudies} (bottom), our results suggest that allowing each gate to be independently learnable and input-dependent yields progressively better outcomes. More detailed analyses of each backbone component can be found in Appendix~\ref{sec:appAblDetails}.

\subsection{xLSTM as Large Language Model}
\label{sec:ExpLanguage}

In this work, we extend our analysis to comprehensive language modeling experiments, focusing on evaluating xLSTM’s effectiveness as a large language model (LLM). To facilitate this, we substantially scale the training data, utilizing 300B tokens from SlimPajama. This token budget aligns with other leading approaches such as Mamba~\citep{Gu:24arxiv} and Griffin~\citep{De:24arxiv}.

Our experimental comparisons feature xLSTM alongside RWKV-4, Llama, and Mamba, which represent distinct method classes as outlined in Section~\ref{sec:ExpComparison}. RWKV-4 is selected as the recurrent neural network benchmark, since appropriate precision training settings for RWKV-5, RWKV-6, and HGRN2 became available only after the commencement of our 300B token runs (further details are provided in Appendix~\ref{sec:appGenTrainProc}).

We implement several model configurations (125M, 350M, 760M, 1.3B), systematically testing each for their ability to extrapolate sequence length, measuring performance on the validation dataset, evaluating downstream task outcomes, and examining language modeling results across the PALOMA benchmark’s 571 text domains. The investigation concludes with an analysis of scaling law patterns exhibited by the models.

\paragraph{Sequence Length Extrapolation.}
\label{sec:spaj300B_ctx_extrapolation}
We begin by evaluating sequence length generalization for large-scale (1.3B parameter) models—specifically xLSTM, RWKV-4, Llama, and Mamba. Each model is trained on sequences with a maximum context of 2048 tokens and then assessed on extended context lengths, reaching up to 16384 tokens. The outcomes of these experiments are visualized in Figure~\ref{fig:spaj300B_seq_extrapolation}. Notably, xLSTM sustains consistently low perplexity scores across these longer sequence lengths, distinguishing itself from the other architectures.

\paragraph{Validation Perplexity and Downstream Tasks.}
\label{sec:spaj_downstream_eval}
Next, across a range of model scales, we systematically compare xLSTM, RWKV-4, Llama, and Mamba in terms of their performance on next-token prediction using the SlimPajama validation data, as well as on a suite of downstream tasks targeting common sense reasoning abilities.

The validation set perplexities for each method are shown in the third column of Table~\ref{tab:lmeval}. Across every model size, xLSTM[1:0] and xLSTM[7:1] consistently achieve the lowest perplexity scores. Performance metrics for various downstream tasks are reported in the remaining columns of Table~\ref{tab:lmeval}. With the exception of occasional superiority by Mamba on the ARC task, xLSTM outperforms other methods in nearly all tasks and model configurations. Additional information can be found in Appendix~\ref{sec:appExpLanguage}.

\paragraph{Performance on PALOMA Language Tasks.} In our third analysis, we evaluate the next token prediction accuracy of xLSTM, RWKV-4, Llama, and Mamba models across different parameter sizes on PALOMA language tasks~\citep{Magnusson:23arxivshort}. The metric used is perplexity, assessed over 571 diverse text domains, ranging from sources like nytimes.com to Reddit's r/depression community.

Results are presented in Table~\ref{tab:ppleval}, which organizes perplexity outcomes into categories for general language modeling (first seven columns) and specific domain benchmarks (last five columns). On these tasks, xLSTM[1:0] achieves superior results compared to xLSTM[7:1].

xLSTM[1:0] exhibits lower perplexity than Mamba on 568 out of 571 domains (99.5\%), outperforms Llama in 486 domains (85.1\%), and achieves better perplexity than RWKV-4 in 570 domains (99.8\%). Further methodological details are provided in Appendix~\ref{sec:appExpLanguage}.

\paragraph{Scaling Laws.} Next, we analyze the power-law scaling phenomena, which enable forecasting model performance as the parameter count increases~\citep{Kaplan:20,Brown:20short}.
As depicted in Figure~\ref{fig:temp_sclaw300B}, the scaling trends are consistent across all architectures, although each displays a distinct offset. RWKV-4 yields the lowest results, with Llama and Mamba following. xLSTM outperforms Mamba by a margin similar to the one Mamba has over Llama. This observed scaling behavior implies that, as model sizes grow, xLSTM is expected to maintain advantageous performance relative to Transformer-based and State-Space architectures.

\textbf{Generation Times and Maximal Throughput.} To evaluate text generation efficiency, we analyze both generation duration and maximal throughput, as depicted in Figure~\ref{fig:inference_time_generation} and Figure~\ref{fig:inference_time_decoding_throughput} (left) for the xLSTM models at the 1.3B parameter level. Comparative benchmarks include equivalently scaled Mamba, Llama, and RWKV variants sourced from HuggingFace, utilizing a static key-value cache for Llama. During experimentation, neither full cache compilation for the Transformer nor Mamba model compilation via \texttt{torch.compile} was functional. All text generation assessments employ a batch size of 1 with a pre-fill of 16 across models—a configuration optimized to benefit Transformer-based models.

The linear time complexity associated with xLSTM, Mamba, and RWKV-4 contrasts distinctly with Llama’s quadratic growth, as illustrated in Figure~\ref{fig:inference_time_generation}. For decoding throughput, multiple batch size and pre-fill combinations were tested on the Llama model. Figure~\ref{fig:inference_time_decoding_throughput} (right) demonstrates that, owing to its fixed memory requirements, xLSTM is capable of supporting substantially larger batch sizes relative to Llama, resulting in superior throughput performance.

\section{Limitations}

(i) Unlike mLSTM, sLSTM’s method of mixing memory inherently prevents operations from being executed in parallel, thus hindering efficient parallel computation. However, we implemented an accelerated CUDA kernel for sLSTM, resulting in a runtime that is less than twice as slow as our parallel mLSTM kernel.

(ii) The mLSTM CUDA kernels lack thorough optimization, which leads to a current performance approximately four times slower than FlashAttention or the scanning mechanism employed in Mamba. Improved CUDA kernels, in line with FlashAttention techniques, may yield considerably better speed.

(iii) mLSTM’s matrix memory architecture entails substantial computational cost due to processing $d \times d$ matrices. Nonetheless, since memory updating and retrieval operate without learned parameters and can leverage standard parallelized matrix routines, the practical impact on wall clock time remains minimal.

(iv) The initialization strategy for the forget gates requires deliberate consideration.

(v) The matrix memory of mLSTM does not scale with sequence length, which presents the risk of memory becoming saturated with longer contexts. Nevertheless, empirical results indicate this is not a concern for context sizes up to 16k; for further discussion, see Section~\ref{sec:spaj300B_ctx_extrapolation}.

(vi) The computational demands for large-scale language modeling experiments prevented comprehensive architecture and hyperparameter optimization, particularly for the larger xLSTM variants. We expect that thorough optimization will be crucial for xLSTM to achieve maximal performance.

\section{Conclusion}
We have made progress in addressing our primary inquiry: What are the capabilities of language modeling when LSTM architectures are scaled to billions of parameters? To date, our findings suggest: ``They can achieve performance comparable to existing approaches such as Transformers and State Space Models''. We extended LSTM into xLSTM through exponential gating combined with memory mixing and introduced a novel memory architecture. Empirical results demonstrate that xLSTM models surpass or match state-of-the-art benchmarks, including Transformers and State Space Models, in language modeling tasks. Scaling patterns observed imply that expanding the size of xLSTM architectures could position them as formidable contenders to prevailing Large Language Models based on Transformer design. Moreover, xLSTM may offer substantial advances in related domains such as Reinforcement Learning, Time Series Forecasting, and the representation of physical systems.

\end{document}